\documentclass[../../main.tex]{subfiles}
\begin{document}
\begin{oframed}
  \subfile{problem}
\end{oframed}

{\bf [解答]} $n$回目に$3$人，$2$人でジャンケンする確率をそれぞれ$a_n$，$b_n$とおく．
また，求める確率を$c_n$とおく．ジャンケンの推移する確率は以下の通り．
\begin{center}
  \begin{table}[htb]
    \centering
    \caption{じゃんけんの推移確率}
    \begin{tabular}{cc}
      3人$\to$3人  & $1/3$\\
      3人$\to$2人  & $1/3$\\
      3人$\to$1人  &$1/3$ \\
      2人$\to$2人  &$1/3$ \\
      2人$\to$1人  &$2/3$
    \end{tabular}
  \end{table}
\end{center}
よって，ジャンケンの推移から以下の漸化式を得る．
\begin{subnumcases}
  {}
  c_n=\dfrac{1}{3}a_n+\dfrac{2}{3}b_n \label{1}\\
  a_{n+1}=\dfrac{1}{3}a_n  \label{2}\\
  b_{n+1}=\dfrac{1}{3}a_n+\dfrac{1}{3}b_n \label{3}
\end{subnumcases}
さらに，初期条件は$a_1=1$，$b_1=0$である．従って\cref{2}から
\begin{align}
  a_n=\left(\dfrac{1}{3}\right)^{n-1} \label{eq:1}
\end{align}
となる．さらにこれを\cref{3}に代入して
\begin{align}
  b_{n+1}&=\dfrac{1}{3}b_n+\left(\dfrac{1}{3}\right)^n
\end{align}
である．$d_n=3^nb_n$とおいて両辺に$3^{n+1}$をかけると
\begin{align}
  d_{n+1}&=d_n+3 \\
  d_n&=3(n-1)+d_1\\
  &=3(n-1)+3b_1 \\
  &=3(n-1)  \tag{$\because b_1=0$}\\
  \therefore
  b_n&=3(n-1)\left(\dfrac{1}{3}\right)^{n}  \\
  &= (n-1)\left(\dfrac{1}{3}\right)^{n-1}
\end{align}
を得る．ただし$d_n=3^nb_n$である．これを\cref{1}に代入すれば求める確率は
\begin{align}
  c_k=(2k-1)\left(\dfrac{1}{3}\right)^k \tag{答}
\end{align}
である．

{\bf[解$2$]}まず$k\ge2$とする．事象$A$，$B$を
\begin{itemize}
  \item $A\cdots$  $3$人でジャンケンする．
  \item $B\cdots$  $2$人でジャンケンする．
\end{itemize}
と定める．$k$回目に一人の勝者が決まるのは以下のいずれかである．
\begin{enumerate}[(i)]
  \item $A\to \dots\to \stackrel{a回目}{A}\to B\to \dots \to \stackrel{k回目}{B}$
  \item $A\to\dots\to \stackrel{k回目}{A}$
\end{enumerate}
(i)で$A$を行う回数$a$，$B$を行う回数を$b$とすると，
\begin{align}
  a+b=k\label{eq:2}
\end{align}
が成り立つ．ただし$a,b>0$である．[解] でのジャンケンの推移の確率から，(i)となる確率は
\begin{align}
  p(a)=\left(\dfrac{1}{3}\right)^{a-1}\dfrac{1}{3}\left(\dfrac{1}{3}\right)^{b-1}\dfrac{2}{3}%
  =2\left(\dfrac{1}{3}\right)^k\tag{$\because\cref{eq:2}$}
\end{align}
同様に(ii)となる確率は
\begin{align}
  q=\left(\dfrac{1}{3}\right)^{k-1}\dfrac{1}{3}=\left(\dfrac{1}{3}\right)^k
\end{align}
となる．故に求める確率は
\begin{align}
  \sum_{a=1}^{k-1}p(a)+q=(2k-1)\left(\dfrac{1}{3}\right)^k \tag{答}
\end{align}
となる．
\newpage
\end{document}
